\subsection{First Theorems on Peano Systems}

The point of introducing Peano systems abstractly is that the first structural
theorems can now be stated once and reused anywhere a Peano system is later
recovered. The first payoff should be the working induction principle that
translates the subset-form axiom into the predicate form used in proofs.

\begin{topicbox}{Induction as a Working Principle}
\begin{exposition}
The induction axiom for a Peano system is stated in terms of subsets of the
underlying set. In actual proofs, however, one usually reasons with a
predicate or property rather than by naming the corresponding subset
explicitly. This topic extracts that predicate-form induction principle from
the subset-form axiom so that the rest of the natural-numbers chapter can use
induction in its ordinary working form.
\end{exposition}

\begin{tcolorbox}[colback=thmbox, colframe=thmborder, arc=2pt,
  left=6pt, right=6pt, top=4pt, bottom=4pt,
  title={\small\textbf{Theorem (Induction Principle for a Peano System)}},
  fonttitle=\small\bfseries]
\begin{theorem}[Induction Principle for a Peano System]
\label{thm:induction-principle-for-peano-system}
Let $(P,S,1)$ be a Peano system, and let $\varphi$ be a predicate on $P$.
If
\[
\varphi(1)
\qquad\text{and}\qquad
\forall n \in P\;(\varphi(n) \Longrightarrow \varphi(S(n))),
\]
then
\[
\forall n \in P\;\varphi(n).
\]
\end{theorem}
\end{tcolorbox}

\begin{remark*}[Standard quantified statement]
\[
\begin{aligned}
&\text{Let } (P,S,1) \text{ be a Peano system and let } \varphi
\text{ be a predicate on } P.\\
&\Bigl(\varphi(1) \land
\forall n \in P\;(\varphi(n) \Longrightarrow \varphi(S(n)))\Bigr)
\Longrightarrow
\forall n \in P\;\varphi(n).
\end{aligned}
\]
\end{remark*}

\begin{remark*}[Contrapositive quantified statement]
\[
\begin{aligned}
&\text{Let } (P,S,1) \text{ be a Peano system and let } \varphi
\text{ be a predicate on } P.\\
&\neg\forall n \in P\;\varphi(n)
\Longrightarrow
\neg\Bigl(\varphi(1) \land
\forall n \in P\;(\varphi(n) \Longrightarrow \varphi(S(n)))\Bigr).
\end{aligned}
\]
\end{remark*}

\begin{remark*}[Interpretation]
This theorem is the proof-ready form of the induction axiom. Instead of
speaking about an inductive subset of the underlying set, it allows one to
argue directly with a property. The mathematical content is the same: if the
property holds at the distinguished first element and is inherited by the
successor operation, then it holds everywhere in the Peano system.
\end{remark*}

\begin{remark*}[Dependencies]
\hyperref[def:peano-system]{Definition~\ref*{def:peano-system}},
\hyperref[ax:peano-induction]{Axiom~\ref*{ax:peano-induction}}.
\end{remark*}

\begin{remark}[Separation of roles]
This theorem belongs here because it is an immediate structural consequence of
the Peano axioms. The later induction chapter should still gather the broader
toolkit: strong induction, induction from an arbitrary starting point,
least-counterexample reasoning, and the equivalence among the standard forms.
\end{remark}
\end{topicbox}
